% A1. Reading Claims and Designing Alternatives: student handout.
% Build: tectonic A1_brief.tex
\documentclass[11pt,a4paper]{article}

\usepackage[margin=2.3cm]{geometry}
\usepackage{graphicx}
\usepackage{booktabs}
\usepackage{tabularx}
\usepackage{array}
\usepackage{enumitem}
\usepackage{xcolor}
\usepackage[colorlinks=true,urlcolor=coral,linkcolor=ink]{hyperref}
\usepackage{fancyhdr}

\definecolor{ink}{HTML}{23272B}
\definecolor{coral}{HTML}{E8543F}
\definecolor{mist}{HTML}{6B7480}
\definecolor{panel}{HTML}{F4F1EA}

\renewcommand{\familydefault}{\sfdefault}
\color{ink}
\setlength{\parindent}{0pt}
\setlength{\parskip}{0.55em}

\usepackage{sectsty}
\allsectionsfont{\sffamily\bfseries\color{ink}}
\sectionfont{\Large\sffamily\bfseries\color{ink}}

\pagestyle{fancy}
\fancyhf{}
\fancyhead[L]{\color{mist}\footnotesize 36104 Data Visualisation and Narratives · Spring 2026}
\fancyhead[R]{\color{mist}\footnotesize A1 · Reading Claims and Designing Alternatives}
\fancyfoot[C]{\color{mist}\footnotesize\thepage}
\renewcommand{\headrulewidth}{0pt}

\newcommand{\labelbox}[1]{\colorbox{panel}{\texttt{\small #1}}}
\newcommand{\poolartefact}[7]{%
  {\large\bfseries #1 · #2}\par
  {\color{mist}\small Domain: #3 · #4 · #5}\par\vspace{0.35em}
  \begin{center}
    \includegraphics[width=0.94\textwidth,height=0.31\textheight,keepaspectratio]{#6}
  \end{center}
  {\footnotesize\textbf{Source:} \href{#7}{Open the chart and source page}\par}
}

\begin{document}

% ---------------- title block ----------------
{\color{coral}\bfseries\small 36104 · ASSESSMENT 1}\\[0.2em]
\resizebox{\textwidth}{!}{\bfseries Reading Claims and Designing Alternatives}\\[0.4em]
{\large\color{mist} See what a chart claims, create an alternative view, and verify a machine's reading of its data.}

\vspace{0.9em}
\begin{tabularx}{\textwidth}{@{}l X@{}}
\toprule
\textbf{Weight} & 30\% of the subject · individual \\
\textbf{Covers} & Sessions 1-3 \\
\textbf{Artefacts} & Three charts from the released pool (twelve entries across varied domains), from three different domains. Artefact 1 is the \textbf{primary artefact} used for the supervised critique. \\
\textbf{Released} & August 8 \\
\textbf{Supervised component} & 25-minute AI-restricted window at the start of the fourth live session (Week 5). Thursday 27 August 2026 \\
\textbf{Due} & \textbf{Friday 28 August 2026, 11:59\,pm} (Sydney time) \\
\textbf{Submit} & One ZIP: \labelbox{A1\_template.ipynb} · the \labelbox{data/} folder · the three supplied chart images renamed \labelbox{original\_chart\_1.png} to \labelbox{original\_chart\_3.png} \\
\bottomrule
\end{tabularx}

\section*{The task}

You will take \textbf{three sourced chart artefacts from the released pool,
from three different domains}, and work them through the course pipeline: \emph{see}
what each claims, develop an alternative view with the Visual Vocabulary, and
\emph{verify} a machine's interpretation of the data. For \textbf{each artefact},
complete a full 9-field critique, an alternative view, a full Chart Choice
Decision Record, and a claim audit. Artefact 1 is the \textbf{primary artefact}:
the one whose critique is completed in the supervised, AI-restricted window.
That supervised critique is the only requirement that differs across the three.

\subsection*{Part A. Seeing (the critiques)}

For each artefact, complete the full 9-field critique:
main claim; audience; visual task; directly visible; interpretation rather
than observation; what is omitted; what misleads; what needs verifying; and
a repaired caption. For artefact 1---the \textbf{primary artefact}---address all
nine fields by hand in the supervised, AI-restricted window on 27 August, then
transcribe and edit them for clarity without AI in \labelbox{critique\_1}.
Complete \labelbox{critique\_2} and \labelbox{critique\_3} outside that window.

\textbf{Critique field guide.} \emph{Main claim}: the one-sentence message the
chart advances. \emph{Audience}: the intended reader or user. \emph{Visual
task}: what the reader must do---for example compare, rank, locate or track
change. \emph{Directly visible}: what is literally encoded in marks, labels and
scales. \emph{Interpretation rather than observation}: what the design invites
you to infer but does not directly show. \emph{What is omitted}: relevant data,
context, uncertainty or comparison left out. \emph{What misleads}: how selection,
encoding or framing---including title, caption, annotation or baseline---could
create an unjustified impression. \emph{What needs verifying}: the source,
calculation or inference that requires checking. \emph{Repaired caption}: a
caption that states the supported finding and its important limit.

\subsection*{Part B. Choosing (the alternative views)}

For each artefact, use the supplied source-data extract and state a claim that
can be supported by that data. Create an alternative view designed to
communicate that claim. Across the three artefacts, the alternatives must use
valid Visual Vocabulary categories and represent \textbf{at least two different
categories}. If you use the same category more than once, explain in each
Decision Record why it is the best fit for that artefact's claim. Explain what each design
makes easier and harder to see than the published original and how those
differences follow from the stated claim. Complete the full \textbf{Chart
Choice Decision Record} for each artefact
(\labelbox{decision\_record\_1..3}).

\textbf{Decision Record field guide.} Identify the intended audience; the
question they need answered; the decision the view supports; the Visual
Vocabulary category; the comparison the reader must make; and the selected
chart form. Then explain why that form fits, name one genuine alternative you
considered, and give a specific reason for rejecting it.

\subsection*{Part C. Verifying (three claim audits)}

Generate an interpretation of each dataset using the \textbf{supplied prompt}
(verbatim, in the template). Classify \textbf{every claim} with the four-way
taxonomy from the Seeing Data lab:

\begin{center}
\begin{tabularx}{0.92\textwidth}{@{}l X@{}}
\toprule
\textbf{Label} & \textbf{Meaning} \\
\midrule
\labelbox{supported} & the data directly backs it, attach an evidence probe \\
\labelbox{plausible but unverified} & sounds right; this dataset cannot settle it \\
\labelbox{unsupported} & asserted with no evidence in this data either way \\
\labelbox{contradicted} & the data shows otherwise, attach an evidence probe \\
\bottomrule
\end{tabularx}
\end{center}

Probes are zero-argument functions returning the evidence (a filtered or
aggregated DataFrame/Series): required for every \labelbox{supported} and
\labelbox{contradicted} claim. Each audit must contain at least one claim that
can be tested with a probe; do not manufacture other claim types simply to
fill a quota. Finish with the five-question AI disclosure.

In the notebook, each claim is a short paraphrase used as the matching key in
both the classification dictionary and, when required, the probe dictionary.

\section*{Your artefacts}

Choose \textbf{three artefacts from the released pool, from three different
domains}. Each of the twelve entries contains a published or source-linked
chart, corresponding tidy source-data extract, data dictionary and provenance
sheet. The approved printable chart-card examples are the only course-reproduction
exception. Use the supplied data rather than digitising the image. Choose your primary before the
supervised session and complete the data preparation, alternative view and other
preparation in advance. On 27 August, record its pool ID on the supervised
sheet. The pool ID entered for artefact 1 in the notebook---for example,
\labelbox{POOL-03}---must match the pool ID written on that sheet. Two students
may study the same source chart, but
every submission must be individual and independently produced.

Use the following canonical domain labels exactly in \labelbox{META\_n}:

\begin{center}
\begin{tabular}{@{}ll ll@{}}
\toprule
\textbf{Pool ID} & \textbf{Domain} & \textbf{Pool ID} & \textbf{Domain} \\
\midrule
POOL-01 & climate science & POOL-07 & economics \\
POOL-02 & public health & POOL-08 & ecology \\
POOL-03 & astronomy & POOL-09 & astronomy \\
POOL-04 & wellbeing & POOL-10 & energy \\
POOL-05 & economics & POOL-11 & geophysics \\
POOL-06 & energy & POOL-12 & transportation \\
\bottomrule
\end{tabular}
\end{center}

\textbf{Known source limitations.} POOL-04's chart card shows ranks 1--48
while the supplied source contains all 143 ranked countries and does not
contain the simultaneous rank intervals drawn in the chart. In POOL-06, a few
published Sankey component labels do not add exactly to the displayed node
totals; treat this as a documented source-chart limitation, not as a student
calculation error. POOL-12 contains undirected historical airport-pair
records: its two endpoints are alphabetically ordered, and
\labelbox{airline\_record\_count} is neither passenger nor flight volume. The
individual provenance sheets repeat these details.

\textbf{Background research and domain knowledge.} The chart, source data and
pool notes are starting points, not a complete package of everything you may
need. Answering some questions well may require modest background reading about
the topic, measures and context represented in the visualisation. You are not
expected to become a domain expert, but you should learn enough to engage
intelligently with the chart's claims, assumptions and limitations. You may use
web search, library resources, reports and other credible sources for this
purpose. If you use an AI system to locate or explain background information,
independently verify its factual claims against credible sources. Keep a record
of the sources that inform your analysis and cite them where relevant.
You should also search for other visualisations of the same or closely related
data to see how other designers have framed, filtered and encoded it and to
develop ideas for your own alternative. Do not copy a design. If another
visualisation materially influences your choices, identify it briefly in the
relevant Decision Record---normally under
\labelbox{why\_it\_fits} or \labelbox{reason\_for\_rejection}---and provide a
working URL. A separate reference list is not required.
Background research should support your own analysis of the chart; it does not
replace close examination of the visualisation and its data.

\textbf{Data-use rule.} Copy the supplied CSV for each chosen chart into the
submission's \labelbox{data/} folder, rename it with its pool ID, and load that
filename in \labelbox{data\_n}. Copy the publisher, publication value and
source URL from the provenance sheet into the corresponding
identifying-information block (\labelbox{META\_n}). Keep \labelbox{data\_n} as
the unchanged DataFrame loaded from that exact file. You may create separately
named filtered, reshaped or derived variables, but must preserve the source
fields and document every transformation. Do not add values inferred
from the image or present derived values as source measurements. Data handling
and disclosure are judged under H2.

\textbf{Reconstruction limit.} The supplied extract may not reproduce every
detail of the published original. Establish which components are supported by
the supplied fields and which depend on unavailable data, transformations,
modelling, annotations or design assets. Use this judgement when critiquing the
original and explaining your alternative.

\textbf{Copyright note.} Include only the chart image needed for criticism and
review, cite its publisher, creator (where known), date and source URL, and do
not republish the assessment beyond the subject site. Follow any additional
UTS library or copyright guidance supplied with the task.

\section*{AI policy}

Part A's primary critique is drafted AI-free in the supervised window and may
be edited for clarity without AI after transcription. An assistant may be
used for Part B's coding and alternative-view work; you remain responsible for
checking and accepting or rejecting its output. Part C requires three
assistant-generated interpretations. Report your use or non-use of AI
accurately. False or misleading disclosure is an academic-integrity issue.

\section*{Submission requirements}

The template notebook contains a set of required named variables and functions
for each artefact. These include the identifying-information blocks
(\labelbox{META\_1}, \labelbox{META\_2}, and \labelbox{META\_3}), critiques
(\labelbox{critique\_1} to \labelbox{critique\_3}), datasets
(\labelbox{data\_1} to \labelbox{data\_3}), alternative views
(\labelbox{redesign\_1} to \labelbox{redesign\_3}), Decision Records, claim
audits, evidence probes, and the AI disclosure. \textbf{Do not rename these
items}, because the autograder checks them directly. The autograder executes
your notebook top-to-bottom on a clean machine. \textbf{Run the self-check cell
before submitting}: it runs the
marker's G1--G6 checks (54 marks). G0's remaining 6 marks are awarded only
after a clean top-to-bottom execution on the marking system. An execution
error loses the G0 marks and may also prevent checks that depend on the
affected code; it does not automatically erase otherwise verifiable marks.
If an execution failure prevents fair assessment, the marker may allow one
technical resubmission within 48 hours of notification. Only changes needed
to restore execution are permitted; the analytical content must remain
unchanged.

\section*{Rubric, 100 marks (worth 30\% of the subject)}

The rubric maximum is the sum of the criterion maxima: 60 automated marks plus
40 human-marked marks equals 100. Canvas records both the assignment and rubric
out of 100. The Assignment 1 group contributes 30\% of the subject, so no
manual score conversion is required.

\subsection*{Automated (60 marks, scored by \texttt{a1\_autograder.py})}

\begin{tabularx}{\textwidth}{@{}l >{\raggedright\arraybackslash}p{4.1cm} c X@{}}
\toprule
\textbf{ID} & \textbf{Item} & \textbf{Marks} & \textbf{Check} \\
\midrule
G0 & Notebook executes end-to-end & 6 & all cells run clean, no errors \\
\addlinespace[0.35em]
G1 & Three artefacts: identity and files valid & 6 & \labelbox{STUDENT\_ID} and all identifying-information blocks complete; three distinct IDs from \labelbox{POOL-01..12}; the corresponding canonical domains are distinct; URLs and \labelbox{data\_status} valid; each supplied original chart matches its selected pool item \\
\addlinespace[0.35em]
G2 & Data provenance $\times$3 & 9 & each exact supplied \labelbox{data/POOL-XX\_source\_data.csv} loads unchanged into \labelbox{data\_n}; $\geq$ 6 rows $\times$ $\geq$ 2 columns; no empty columns; source identifying information and transformations documented \\
\addlinespace[0.35em]
G3 & Alternative-view mechanics $\times$3 & 15 & each alternative view (\labelbox{redesign\_n}) returns its Axes; all categories are valid and at least two different Vocabulary categories are represented across the three artefacts; bar-family charts include a zero baseline; every chart title has at least 15 characters \\
\addlinespace[0.35em]
G4 & Claim-audit structure $\times3$ & 14 & for each artefact: a 150--250-word interpretation from the supplied prompt; $\geq$ 4 validly labelled claims; at least one \labelbox{supported} or \labelbox{contradicted} claim; an executable, non-empty probe for every such claim \\
\addlinespace[0.35em]
G5 & Critiques + Decision Records $\times3$ & 6 & all three 9-field critiques and all three 9-field Decision Records are complete; each \labelbox{vocabulary\_category} matches \labelbox{CATEGORY\_n} \\
\addlinespace[0.35em]
G6 & Five-question AI disclosure & 4 & all 5 answers completed; each answer has at least 3 words and is not a placeholder \\
\midrule
 & \textbf{Total} & \textbf{60} & \\
\bottomrule
\end{tabularx}

\subsection*{Human-marked (40 marks, anchored 0/1/2 per item)}

\begin{tabularx}{\textwidth}{@{}>{\raggedright\arraybackslash}p{2.9cm} >{\raggedright\arraybackslash}X >{\raggedright\arraybackslash}X >{\raggedright\arraybackslash}X@{}}
\toprule
\textbf{Item (weight)} & \textbf{0} & \textbf{1} & \textbf{2} \\
\midrule
H1 · Critique insight across three ($\times$5) & Restates the charts & Correctly identifies claims and real omissions & Also separates \emph{directly visible} from \emph{interpretation}, explains how title, caption, annotation or baseline frames the claim where relevant, and the repaired captions actually repair \\
\addlinespace[0.45em]
H2 · Data handling, across artefacts ($\times$2.5) & Source data misread or transformations undisclosed & Units and transformations mostly correct & Source fields preserved; transformations correct, reproducible and disclosed \\
\addlinespace[0.45em]
H3 · Decision Records across three ($\times$5) & No argument, or category labels only & Correct task categories with generic justifications & Ties audience $\rightarrow$ task $\rightarrow$ channel precision; rejected alternatives are real and reasons specific \\
\addlinespace[0.45em]
H4 · Alternative-view craft, across the three ($\times$5) & Unreadable or misleading & Honest defaults, readable & Titles state findings; scales honest; labels carry the argument \\
\addlinespace[0.45em]
H5 · Audit judgement across three ($\times$2.5) & Classifications mostly wrong & Labels mostly appropriate, but probes are weak & Labels distinguish evidence from plausibility precisely; every required probe directly tests the claim \\
\midrule
\multicolumn{4}{@{}l}{\textbf{Total: 40 marks} (weights $\times$ anchor score; maximum $2\times(5+2.5+5+5+2.5)=40$)} \\
\bottomrule
\end{tabularx}

\subsection*{Process requirements (not marked, but gates)}

\begin{itemize}[leftmargin=1.4em,itemsep=0.15em]
  \item Part A supervised notes handed in during the 27 August window.
        Missing notes $\rightarrow$ H1 capped at 0.
  \item The pool ID entered for artefact 1 in the notebook must match the pool
        ID written on the supervised sheet (27 August).
\end{itemize}

\clearpage
\section*{Artefact pool: charts and sources}

The twelve chart images are reproduced below for assignment reference. The
downloadable pool supplies the corresponding chart image, CSV, data dictionary
and provenance sheet for each entry.

\poolartefact{POOL-01}{Deviation}{climate science}{Global mean temperature anomaly}{NASA Goddard Institute for Space Studies}{pool/POOL-01/original_chart.png}{https://data.giss.nasa.gov/gistemp/graphs_v4/}
{\footnotesize\textbf{Data:} \href{https://data.giss.nasa.gov/gistemp/graphs_v4/graph_data/Global_Mean_Estimates_based_on_Land_and_Ocean_Data/graph.csv}{Download the publisher's source data}\par}
\vfill
\poolartefact{POOL-02}{Correlation}{public health}{Life expectancy vs health expenditure}{Our World in Data}{pool/POOL-02/original_chart.png}{https://ourworldindata.org/grapher/life-expectancy-vs-health-expenditure}
{\footnotesize\textbf{Data:} \href{https://ourworldindata.org/grapher/life-expectancy-vs-health-expenditure.csv}{Download the publisher's source data}\par}

\clearpage
\poolartefact{POOL-03}{Correlation}{astronomy}{Hubble's 1929 distance--velocity relation}{Proceedings of the National Academy of Sciences}{pool/POOL-03/original_chart.png}{https://pmc.ncbi.nlm.nih.gov/articles/PMC522427/}
{\footnotesize\textbf{Data:} \href{https://apod.nasa.gov/diamond_jubilee/1996/hub_1929.html}{Open the published tabulated observations}\par}
\vfill
\poolartefact{POOL-04}{Ranking}{wellbeing}{Country rankings by life evaluation, 2021--2023}{World Happiness Report}{pool/POOL-04/original_chart.png}{https://worldhappiness.report/ed/2024/}
{\footnotesize\textbf{Data:} \href{https://files.worldhappiness.report/WHR24_Data_Figure_2.1.xls}{Download the report's supporting data}\par}

\clearpage
\poolartefact{POOL-05}{Distribution}{economics}{Lorenz curves for US family income}{OpenStax}{pool/POOL-05/original_chart.png}{https://openstax.org/books/principles-economics-3e/pages/15-4-income-inequality-measurement-and-causes}
{\footnotesize\textbf{Data:} \href{https://openstax.org/books/principles-economics-3e/pages/15-4-income-inequality-measurement-and-causes}{Open the publisher's chart and source table}\par}
\vfill
\poolartefact{POOL-06}{Flow}{energy}{Estimated US energy consumption, 2023}{Lawrence Livermore National Laboratory / US Department of Energy}{pool/POOL-06/original_chart.png}{https://flowcharts.llnl.gov/commodities/energy}
{\footnotesize\textbf{Data:} \href{https://flowcharts.llnl.gov/commodities/energy}{Open the published diagram; flow values are printed on the chart}\par}

\clearpage
\poolartefact{POOL-07}{Change over time}{economics}{Measures of underlying inflation}{Reserve Bank of Australia}{pool/POOL-07/original_chart.png}{https://www.rba.gov.au/chart-pack/aus-inflation.html}
{\footnotesize\textbf{Data:} \href{https://www.rba.gov.au/statistics/tables/csv/g1-data.csv}{Download RBA table G1 data}\par}
\vfill
\poolartefact{POOL-08}{Correlation}{ecology}{Lynx--hare phase portrait}{San Diego State University (J. M. Mahaffy)}{pool/POOL-08/original_chart.png}{https://jmahaffy.sdsu.edu/courses/f09/math636/lectures/lotka/qualde2.html}
{\footnotesize\textbf{Data:} \href{https://jmahaffy.sdsu.edu/courses/f09/math636/lectures/lotka/qualde2.html}{Open the published observations}\par}

\clearpage
\poolartefact{POOL-09}{Magnitude}{astronomy}{Cosmic abundance of the elements}{NASA Imagine the Universe / HEASARC}{pool/POOL-09/original_chart.png}{https://imagine.gsfc.nasa.gov/educators/elements/imagine/cosmic_abundances.html}
{\footnotesize\textbf{Data:} \href{https://chem.libretexts.org/Bookshelves/General_Chemistry/Interactive_Chemistry_(Moore_Zhou_and_Garand)/06\%3A_Appendix/6.02\%3A_Elemental_Abundances}{Open the referenced abundance table}\par}
\vfill
\poolartefact{POOL-10}{Part-to-whole}{energy}{Sources of US electricity generation}{US Energy Information Administration}{pool/POOL-10/original_chart.png}{https://www.eia.gov/energyexplained/electricity/electricity-in-the-us.php}
{\footnotesize\textbf{Data:} \href{https://www.eia.gov/energyexplained/electricity/electricity-in-the-us.php}{Open the publisher's chart and source values}\par}

\clearpage
\poolartefact{POOL-11}{Spatial}{geophysics}{Seismicity of the Earth 1900--2007}{US Geological Survey}{pool/POOL-11/original_chart.png}{https://pubs.usgs.gov/sim/3064/}
{\footnotesize\textbf{Data:} \href{https://earthquake.usgs.gov/fdsnws/event/1/query?format=csv\&starttime=1900-01-01\&endtime=2008-01-01\&minmagnitude=7.3\&orderby=time}{Download the documented USGS catalogue extract (1900--2007, M\,$\geq$\,7.3)}\par}
\vfill
\poolartefact{POOL-12}{Flow}{transportation}{Worldwide airline route network}{OpenFlights}{pool/POOL-12/original_chart.png}{https://openflights.org/data.php}
{\footnotesize\textbf{Data:} \href{https://raw.githubusercontent.com/jpatokal/openflights/master/data/routes.dat}{Download the OpenFlights route data}\par}

\end{document}
